\emph{Setting.} $A$ and its generators are those of Theorem \claimref{tower.algebraic-limit-even-charge} and Proposition \claimref{tower.ladder-pairs-satisfy-car-vectors}; states are those of Premise \claimref{synthesis.premise-positivity-of-values}. A vector state of the Fock module (Theorem \claimref{tower.fock-module-one-qubit-per-rung}) satisfies (P). For $x$ involving only the first $n$ rungs, $\langle s_{\le m}, x\,s_{\le m}\rangle$ is the same for all $m \ge n$, since $x$ does not act on the later rungs; so $\omega_{s}$ is well defined, and it is a state.

\emph{(i).} The sign change of a set of generators preserves the defining relations, so it extends to an automorphism of each floor and of $A$, and it commutes with $\dagger$, which is defined generator by generator. The sign change of $f_{j}$ exchanges $v_{j}$ and $w_{j}$, hence $\mathrm{on}_{j}$ and $\mathrm{off}_{j}$, and fixes the other rungs' tests: it carries $\omega_{s}$ to the string state with the value at $j$ changed. Applying it at every rung where $s$ and $t$ differ, a finite set for $G_{1}$ and an arbitrary set for $G_{3}$, carries $\omega_{s}$ to $\omega_{t}$. For $G_{2}$: $e_{j}$ is an odd unitary with $e_{j}v_{j}e_{j} = w_{j}$, and it commutes with the tests of the other rungs, so conjugation by a product of $e_{j}$ over a finite set changes the values on that set. Conversely, an element of $G_{1}$ fixes every generator beyond some floor; and a unitary $u$ of the floor of the first $n$ rungs commutes with $\mathrm{on}_{j}$ for $j > n$, because $\mathrm{on}_{j}$ is an even element involving only rung $j$. In both cases the image of $\omega_{s}$ agrees with $\omega_{s}$ on the tests of all rungs beyond a floor, so if it is a string state $\omega_{t}$ then $t$ agrees with $s$ beyond that floor. Each tail class is countable and there are uncountably many strings. The step's symmetries in Theorem \claimref{synthesis.state-of-a-new-rung-unique-symmetric-extension-is-the-half-weight} are the sign changes of the two generators of one rung; $G_{3}$ is generated by these over all rungs, with infinitely many applied at once allowed. The script checks the finite statements on three rungs.

\emph{(ii).} $x$ is a combination of products of the $v_{i}$, $w_{i}$ of the rungs it involves; these do not change the occupation of any other rung, and occupation vectors with different occupations of a rung are orthogonal. The script verifies this for a full basis of the algebra on two rungs against all pairs of strings on four rungs. For the completed algebra: $\omega_{s}$ is the gauge-invariant quasi-free state whose operator is the diagonal projection with entries $s_{j}$, pure by clause (d) of Fact \claimref{synthesis.imported-powers-stormer-takesaki-disjointness}; the difference of two such projections is diagonal with entries $0, \pm 1$ and is Hilbert-Schmidt exactly when it has finitely many nonzero entries; clause (a) then gives quasi-equivalence exactly for strings in one tail class, and disjointness otherwise. This applies the import and is not proved here.

\emph{(iii).} On a floor of $n$ rungs the invariant extension, iterated, is $x \mapsto \langle x\rangle_{0}$, whose density is $2^{-n}$ times $1$. A family of pairwise orthogonal idempotents with sum $1$, each with scalar part $2^{-n}$, has uniform average with density $2^{-n}\sum = 2^{-n}$: the same state. The $2^{n}$ string projectors are one such family. On two rungs the script verifies that the four stated elements are even, self-adjoint, idempotent, pairwise orthogonal, of scalar part $\tfrac14$ and sum $1$, that each gives two strings the weight $\tfrac12$, and that $\sum_{s}P_{s}\,\cdot\,P_{s}$ sends each to half the sum of two string projectors. The statement about the store is the last sentence of Proposition \claimref{singularity.p-reduced-metric-level-reading-many}: the reading that makes the joint-diagonal algebra the metric level is a choice between two assemblies, not decided there.

\emph{(iv).} With the data a string (or any fixed vector of a finite floor's module), $\omega_{\psi}(y) = \langle\psi \otimes \mathrm{data},\, y\,\psi \otimes \mathrm{data}\rangle$ on a floor containing the rungs $y$ involves, which is $a^{2}M_{00} + 2abM_{01} + b^{2}M_{11}$ with $M_{ij}$ the symmetrised values of $y$ between $\phi_{i} \otimes \mathrm{data}$ and $\phi_{j} \otimes \mathrm{data}$. Suppose $q \in \{0, 1\}$ at every rational point of the unit circle. There are infinitely many such points with pairwise different directions (from the Pythagorean triples), so one of the two values, $\nu$, is taken at infinitely many directions. On the unit circle $q - \nu = q - \nu(a^{2} + b^{2})$, a binary quadratic form that vanishes in infinitely many directions; a nonzero binary quadratic form vanishes in at most two. So $q = \nu(a^{2} + b^{2})$, constant on the circle. No continuity is used. For an automorphism $\alpha$ of $A$ and a test $E$, the weight of $E$ after $\alpha$ is $\omega_{\psi}(\alpha(E))$ with $\alpha(E) \in A$, so it is covered; a pointwise limit of quadratic forms in two variables is a quadratic form, since its three coefficients are determined by the values at three directions. Now suppose a rule assigns, for each value of the data, a definite outcome through such a value, and the data are distributed independently of $\psi$. For each value of the data the outcome is the same for all $\psi$, so the frequency of the first outcome does not depend on $\psi$; the weights of (P) for the two alternatives are $a^{2}$ and $b^{2}$, which do. The script computes, for the four strings of a fresh pair, four couplings and five tests, the eighty forms, confirms each is the weight computed directly at seven rational directions, and that every form that is not a multiple of $a^{2} + b^{2}$ takes a value in $\{0, 1\}$ at no more than four of them (in fact at no more than two).

\emph{(v).} The first index at which $b$ and $d$ differ is $k$ with probability $2^{-k}$, and it is then decided for the first outcome exactly when $d_{k} = 1$; so the probability is $\sum_{k}d_{k}2^{-k} = p$. The index exceeds $k$ with probability $2^{-k}$, and $\sum_{k \ge 0}2^{-k} = 2$. A function of $n$ unbiased values has probabilities with denominator $2^{n}$. The script computes the twenty digits of $\tfrac{9}{25}$, confirms the period and the sum, enumerates all $4096$ prefixes of length $12$ exactly, and confirms that $2^{k}\cdot\tfrac{9}{25}$ is not an integer for $k \le 40$. That the unbiased values could be the values of unread rungs is the reading of clause (iv) of Theorem \claimref{synthesis.selector-candidates-on-the-tower-integer-states-parity-deterministic-laws}; that the tower has no clock is Proposition \claimref{tower.answers-iteration-question}.
